\section{はじめに}
\label{sec:intro}

我々はオープンソース脆弱性スキャナ Vuls\cite{vuls} と，その脆弱性データベース群を継続的に開発・運用している．
前稿\cite{nakaoka2025}では，一次データソース（ベンダやプロジェクトが自ら公開するセキュリティアドバイザリ）を可能な限り直接収集し，取得したままの \texttt{vuls-data-raw} 形式と共通スキーマの \texttt{vuls-data-extracted} 形式に分けてGitで履歴管理する「データハーベスト」方式を提案した\cite{vuls-data-update,vuls-data-db}．
本稿はその基盤の上で，「検知条件を誰から得るべきか」という問いを定量的に扱う．

脆弱性スキャナが判定に用いる情報を本稿では\textbf{検知条件}と呼ぶ．
検知条件とは最低限「どの製品の」「どのバージョン範囲が」影響を受けるかの組であり，これが機械可読な形で与えられなければスキャナはインベントリと突き合わせられない．
NVD\cite{nist-nvd}においてこの役割を担うのが，\texttt{configurations} に格納されるCPE適用条件である．
各条件は \texttt{criteria} にCPE 2.3形式の文字列を持ち，任意で \texttt{versionStartIncluding} / \texttt{versionEndExcluding} 等の範囲指定を伴う．
ただしバージョン部が \texttt{*} や \texttt{-} で範囲指定を持たない場合，その条件は製品全体を影響対象とみなすことしか表現できない．

近年，CVEの報告件数は急増している．
我々の観測では2026年は1月1日から7月29日までの210日間で44,809件（取下げを除く）が公開されており，年率換算で約7.8万件，2023年の28,817件の約2.7倍のペースである．
この増加が解析体制を圧迫していることは，NVDが \texttt{Deferred} というステータスの運用を始めたこと\cite{nvd-status}に表れている．これは遅延ではなくNVDがエンリッチメントを行わないと明示した状態である．
またこの遅延を受け，VulnCheckはNVDのデータに独自の解析・エンリッチメントを加えたVulnCheck NVD++\cite{vulncheck-nvd2}の提供を開始した．NVD由来の \texttt{configurations} を保持したまま独自の結果を \texttt{vcConfigurations} と \texttt{vcVulnerableCPEs} に持つため，両者の寄与を機械的に切り分けられる．

我々のような検知ツールの提供者にとってこれは，「これまで信頼して使えていたデータソースを今後も同じ前提で使えるのか」という運用上の問題である．
しかし我々の知る限り，CVE増加後のNVDおよび商用エンリッチメントについて，検知条件の供給という観点から定量的に整理した報告はない．

そこで本稿は，NVDとVulnCheck NVD++の全件（2026年7月29日時点，取下げを除く353,704件）についてCPE適用条件の付与率と粒度を測り，さらに一次データソースを基準とする4基準を定義してFortinetとCiscoで実測する（\ref{sec:supply}章）．
そのうえで，LLMの非決定性を日々のデータ経路から追い出したまま一次データソースを構造化する3つの適用パターンを整理して実装し（\ref{sec:llm}章），入力データソースのみを変えた4種のデータベースで計10環境を検知して，上記の差が検知結果に現れることを確認する（\ref{sec:eval}章）．
NVDのみでの未検知率はWindows Server 2022の1.3\%からサポート終了版FortiOSの68.8\%まで開き，\ref{sec:supply}章で測った付与率はこの差をほとんど予測しない．
